\section{Open Problems}
\label{sec:open}

% 本节对应 slides/outline.md 的 P11--P13。
% 硬性要求: 2--3 个。每个必须有文献支撑，且必须回答"如果我错了会怎样"。
%
% ⚠️ 回答不了自己提出的争议，比不提出更糟。若确实答不上来，
%    坦白说"这正是我们认为的开放问题"。

\subsection{Q1: Is the reported gain real?}

\textbf{The challenge.} \cite{lin2026notbenefit} argues that harness updating
and harness benefit must be disentangled: a system that can modify its own
harness is not thereby a system that becomes better. We find this position
convincing on inspection of the evidence, for three reasons.

First, a search procedure that generates many candidate harnesses and retains
the best-scoring one will exhibit an apparent gain even if modification is
uninformative; the expected maximum of $N$ draws exceeds their mean.

Second, where the same task set is used both to select among candidates and to
report the result, the reported gain is optimistic by an amount that depends on
the number of candidates.

Third, contamination is difficult to exclude when the modification procedure
itself is given access to benchmark tasks.

\textbf{Our position.} \todo{}

\textbf{What would settle it.} A held-out task set never exposed to the
modification procedure, plus a report of the full distribution of candidate
scores rather than only the maximum.

\textbf{If we are wrong.} \todo{Required. If the gains reproduce
independently, the field's core narrative holds and this review's contribution
reverts to mapping the method space.}

\subsection{Q2: Is the verifier independent of what it verifies?}

\textbf{The challenge.} A harness that both proposes and certifies its own
modifications does not possess an independent check. This is the structural
concern behind axis B's distinction between \code{self-preference} and
\code{external-verifier}.

\cite{pan2026rho} makes the problem concrete by proceeding on the premise that
ground-truth validation is unavailable in deployment, and optimizing the harness
from historical trajectories under self-preference. \cite{banerjee2026severa}
takes the opposite route, attaching formal guarantees to the synthesis of
self-evolving agent programs.

The interesting question is not which route is correct but what each costs.
Self-preference removes a dependency the field may not be able to satisfy in
practice; formal guarantees constrain the space of systems that can be
certified. \todo{After deep-reading SEVerA, state the scope of its guarantee
and its relaxation assumptions.}

\textbf{Our position.} \todo{}

\textbf{What would settle it.} \todo{Has any work re-evaluated self-assessed
results with an independent evaluator?}

\textbf{If we are wrong.} \todo{Required. If self-preference validation works,
the field can operate in entirely unlabelled deployment settings, which is a
major practical advance.}

\subsection{Q3: Is compute accounted for?}

\textbf{The challenge.} \todo{To be supported from the deep-read corpus. This
is the easiest of the three to verify and should be done first.}

\textbf{Our position.} \todo{}

\textbf{If we are wrong.} \todo{Required.}

\subsection{Scope note}

We deliberately omit one candidate concern --- the general reproducibility
crisis of the field, arising from the failure to release complete harnesses and
evolution archives. It is real but engineering-level, and including it would
dilute the argument.

% ----------------------------------------------------------------------
% 候选引用:
%   \cite{lin2026notbenefit}  Q1 的直接出处
%   \cite{pan2026rho}         Q2
%   \cite{banerjee2026severa} Q2 --- 形式化路线
%   \cite{chen2026safety}     安全面
% ----------------------------------------------------------------------
